\section{\uppercase{Introduction}}
\label{sec:intro}

Flying insects are a major sanitary and economic threat in food-industry facilities, where regulatory compliance requires continuous monitoring of fly populations. Adhesive (sticky) traps are the standard monitoring instrument: cheap, passive, and ubiquitous. Their bottleneck is the manual inspection of each plate by technicians who must not only \emph{count} the insects but \emph{identify their species}, since different species imply different contamination sources. Automating this reading is therefore a species-level recognition problem, not merely a counting one.

Computationally, sticky-trap reading sits at the intersection of two regimes that are hard in isolation and rarely studied together. First, it is a \emph{tiny object} problem: in our industrial setting, plates are photographed at 1920$\times$1080 under a fixed-focus protocol, and the mean captured insect measures $\approx$44\,px ($\sqrt{wh}$), about 2\% of the image width, and a large share of the insects fall below the 32\,px limit commonly used to define ``tiny'' objects~\cite{cheng2023soda}; at native resolution, however, each insect still spans several feature-map cells. Second, it is a \emph{fine-grained} problem: the target species differ in subtle cues (wing venation, thorax patterning, size ratios) that degrade quickly as trapped specimens decay~\cite{bockmann2021}. Recent species-level work on sticky traps remains confined to laboratory conditions and single-stage detectors~\cite{laekeman2025}, while the tiny-object literature validates its methods almost exclusively on aerial imagery~\cite{wang2021nwd,xu2022rfla}, and detect-then-classify cascades for species identification are established for vertebrate camera traps~\cite{beery2019} but essentially untested on trap-caught insects.

Prior work introduced a data-centric pipeline for this exact setting~\cite{pacbb2026}: through iterative dataset curation, an explicit \emph{Noise} class and a fixed focal distance, a single-stage YOLO detector was driven from 0.476 to 0.608 mAP. That study concluded, however, that the remaining confusion between visually similar fly species was \emph{irreducible}: no further data curation reduced the misclassification between look-alike pairs. This claim frames the question we address here: \textbf{what limits species-level reading of sticky traps, and which lever moves that limit---data, the identification stage, or the treatment of scale?}

We study three levers under one protocol (same dataset lineage, test split and species-level metrics, multiple seeds): \emph{data} (targeted physical collections), \emph{identification} (a detect-then-classify cascade with five classifiers on four backbones, three of them foundation models, and a ground-truth-crop oracle that bounds any classifier) and \emph{scale} (tiny-object methods from the aerial literature). The answer is mixed and we treat it as a result. Along five cumulative training sets, species accuracy rises from 56\% to 82\%, most visibly where the added collections were aimed; because size, quality and domain co-vary and the equal-size control cannot separate label quality from quantity, we report this as an association, not as proof that any one collection is the cause. Look-alike confusion falls but does not vanish. Identification is not moved: the cascade does not improve on the detector in any of 15 configurations and the oracle saturates near 83\% for every backbone; the residual error is confident, concentrated in the housefly/blowfly/flesh-fly family and shared by all classifiers, which is consistent with limited information in the images (but equally with label noise, the crop recipe or the 44\,px object size, which we cannot separate). None of the aerial tiny-object variants we tried lifts the baseline.

Concretely, we investigate three research questions:

\begin{itemize}
    \item \textbf{RQ1 (data):} Does extending the dataset with targeted physical collections (species-pure, crowded, insect-only plates, white-paper photographs, copy-paste) improve species-level detection and identification? We use a ladder of five cumulative training sets.
    \item \textbf{RQ2 (identification):} Does a detect-then-classify cascade improve on the single-stage baseline, and how far can any classifier go given perfect crops?
    \item \textbf{RQ3 (transfer):} Do tiny-object methods validated on aerial benchmarks transfer to trap-caught insects with a modern end-to-end detector?
\end{itemize}

Our contributions are:
\begin{enumerate}
    \item A data ladder over five cumulative training sets (three seeds each, fixed test split) along which species accuracy rises from 56.1\% to 82.2\%, with look-alike confusion reduced but not removed; an equal-size control shows that random subsamples of the final data match L1 and L2, from L1 to L2 the rise is mostly quantity, but not for L0 (a random 103-image subsample is 19\,pp better).
    \item A controlled comparison of single-stage detection versus a detect-then-classify cascade for species-level insect monitoring on real industrial sticky traps, over five classifier backbones and three detectors, with paired bootstrap intervals over test images. For this recipe (frozen detector, 75\% crop margin, label-only fusion), the cascade does not improve on the detector in any configuration (macro-F1 change $-1.1$ to $-5.1$\,pp on TEST) at 30--67$\times$ the inference time.
    \item A diagnosis of the identification ceiling: ground-truth-crop accuracy saturates at 82.7--83.8\% for all backbones, 4.9\% of instances are misclassified by all fifteen oracle classifiers (five types, three seeds each), 69\% of errors are confident and 79\% fall within the MD/MV/MC family. Changing the classifier does not remove this confusion (which~\cite{pacbb2026} called irreducible, on different data and classes); we state what we can and cannot say about its cause. Part of it is a data effect (MV$\to$MC from 14.8\% to 5.4\% among localized objects), while MC$\to$MV stays near 12\%.
    \item A negative transfer result for default-hyperparameter variants, with three seeds each (one for YOLO11m): eight tiny-object variants of an end-to-end detector (NWD in the loss only or in the assigner only, NWD with RFLA, GCD, Wise-IoU v3, a P2 head, 2560-px input and YOLO11m), tested against a same-campaign control, bring no detectable gain, and NWD with RFLA is 8.8\,pp worse because of unstable seeds; for objects below 32\,px no variant raises species recall. Sliced inference, TTA and ensembling do not beat, or marginally beat ($+0.4$ mAP@50), the native-resolution detector.
\end{enumerate}
